\chapter{Discussion and Future Work}

In this thesis, we propose extending an existing 2D cell segmentation network (Mediar) to 3D by applying an existing slice accumulation technique. We explore ways to improve the performance in both network design and pretraining. By integrating the SAM2 image encoder registry into the Mediar3D network and expanding the pretraining by a large size, we obtain 3 candidate models to benchmark against the baseline, CellposeSAM. We find that our models perform similarly to CellposeSAM, matching or outperforming it while reducing inference time by a factor of 10. Our best-performing model achieved SEG scores of 0.815 and 0.785 on the challenging Fluo-N3DL-TRIC and Fluo-N3DL-DRO datasets of the Cell Tracking Challenge, respectively. While the results look promising, there are some important concerns remaining: First, we only consider the peak performance of our networks, that rely on manually picking a post-processing threshold parameter. This post-processing threshold dramatically influences performance. Although recurring trends can be noticed, such as high thresholds around 0.1–0.5 seem to improve zero-shot performance for highly pretrained models, a specific threshold can not be determined, as the dependence shifts after finetuning iterations. Considering we aim to develop a robust segmentation model, this dependency gives CellposeSAM the upper hand, as their performance seems to be invariant to the probability threshold. Second, although some of our models perform better on the datasets we benchmark on, the test sizes are still too low, in order to draw definite conclusions. Limiting factors for broader benchmarking sizes include data availability and computational costs. Nevertheless, apart from the above mentioned considerations about actual robustness and benchmark validity, we believe our proposed models hold great potential for generalizable, robust cell segmentation networks that can be adapted to several pipelines and downstream tasks, regarding the faster runtimes and promising performance shown in our experiments. We also managed to overcome memory constraints by rewriting the slicing inference loop. For substantially large images (larger than 1000×1000×1000 pixels), post-processing remains to be the bottleneck, as consumer-grade GPUs run out of memory at this point. A reasonable alternative would be, to adapt the learning objective to a tanh formulation, as proposed by  Eschweiler et. al. in \cite{eschweiler2022robust3dcellsegmentation}. Note that, like shifted window post-processing, this method is a trade-off between runtime and performance. 
We have also shown a working pipeline to cope with partially annotated training data by manually removing annotations from a fully annotated dataset. The proposed masked training loss significantly improves training compared to naively using the default pipeline. In practice, success of this method highly depends on the fraction of annotated instances and the target statistics.
For future work, further investigation of threshold dependence on the trained statistics should be made. This way, an automated threshold picking could be established. Since pretraining with larger batch sizes requires either multi-gpu training or gradient accumulation, tuning of sampling ratios is very expensive. Nevertheless, we believe that further studys has the potential to further improve generalization capabilities and potientially robustness towards the cellprobability threshold at test time. 
For our 3D Feature Network, there is a lot of room for improvement. First, slicing for non-axial directions is not introduced yet. By merging slice embeddings with sagittal and coronal planes, the 3D decoder could benefit from more context. Also, from a pretraining point of view, we did not give the network a chance to generalize, as we did not prepare a diverse pretraining dataset for the network. Although annotated 3D data is limited, diffusion models could leverage high annotation costs.